Die Luftwiderstandskraft $F_\text{W} = \frac{1}{2}\,c_w\,\rho\,A\,v^2$ ist proportional zum Quadrat der Geschwindigkeit; $c_w$, $\rho$ und $A$ bleiben gleich. Daher ist
\begin{align}
 \frac{F_2}{F_1} &= \left(\frac{v_2}{v_1}\right)^2
\end{align}
und die relative Zunahme
\begin{align}
 \frac{F_2-F_1}{F_1} &= \pLwF \\
   &= \left(\frac{\vZwO}{\vEiO}\right)^2-1 \\
   &= \pLw \approx \result{\pLwP{0}{2}}
\end{align}
Die Geschwindigkeit nimmt um rund einen Viertel zu, der Luftwiderstand aber um fast zwei Drittel.
